% from ega3-6-fr.tex
\oldpage[III]{146}

est aussi un homomorphisme bijectif. Comme $A'$ est un $A$-module plat, la
conclusion résulte de (6.3.9), et on a de même des isomorphismes canoniques
pour les suites spectrales.
\end{proof}

\subsection{Foncteurs hypertor locaux de complexes de Modules quasi-cohérents : cas des schémas affines}

\begin{env}[6.4.1]
Soient $S$ un schéma affine d'anneau $A$, $X$, $Y$ deux $S$-schémas affines
d'anneaux $B$, $C$ respectivement, de sorte que $B$ et $C$ sont des algèbres
sur $A$. Tout complexe $\mathscr P_\bullet$ (resp. $\mathscr Q_\bullet$) de
$\mathscr O_X$-Modules (resp. $\mathscr O_Y$-Modules) quasi-cohérents est de
la forme $\widetilde{P_\bullet}$ (resp. $\widetilde{Q_\bullet}$), où
$P_\bullet$ (resp. $Q_\bullet$) est un complexe de $B$-modules (resp.
$C$-modules) (I, 1.3.7 et 1.3.8). On peut évidemment considérer $P_\bullet$
et $Q_\bullet$ comme des complexes de $A$-modules et former les
$\mathbf{Tor}_n^A(P_\bullet,Q_\bullet)$; en outre, en vertu du caractère
bifonctoriel de $\mathbf{Tor}_n^A(P_\bullet,Q_\bullet)$, les $A$-algèbres
$B$ et $C$ opèrent dans ce $A$-module, et ces opérations en font un
$(B,C)$-bimodule, ou, ce qui revient au même, un module sur
$B\otimes_A C=A(X\times_S Y)$. On a donc défini de la sorte un
$\mathscr O_{X\times_S Y}$-Module quasi-cohérent
\[
  \mathscr{T}\!or_n^{\mathscr O_S}(\mathscr P_\bullet,\mathscr Q_\bullet)
  =\bigl(\mathbf{Tor}_n^A(P_\bullet,Q_\bullet)\bigr)^{\sim}
  \tag{6.4.1.1}\label{III.6.4.1.1.fr}
\]
que l'on appelle l'\emph{hypertor local d'indice $n$} des complexes
$\mathscr P_\bullet$ et $\mathscr Q_\bullet$ et que l'on note aussi
$\mathscr{T}\!or_n^S(\mathscr P_\bullet,\mathscr Q_\bullet)$.
\end{env}

\begin{namedtheorem}[6.4.2]{Lemme}
Avec les notations de (6.4.1), supposons que l'anneau $A$ soit de la forme
$R^{-1}A'$, où $A'$ est un anneau et $R$ une partie multiplicative de $A'$.
Soit $S'=\operatorname{Spec}(A')$, de sorte que $X$ et $Y$ peuvent être
considérés comme des $S'$-préschémas et que l'on a
$X\times_{S'}Y=X\times_S Y$ (I, 1.6.2 et 3.2.4). On a alors
\[
  \mathscr{T}\!or_\bullet^{S'}(\mathscr P_\bullet,\mathscr Q_\bullet)
  =\mathscr{T}\!or_\bullet^S(\mathscr P_\bullet,\mathscr Q_\bullet).
\]
\end{namedtheorem}

\begin{proof}
Cela résulte de la formule (6.4.1.1) et de (6.3.13).
\end{proof}

\begin{env}[6.4.3]
Avec les notations et hypothèses de (6.4.1), soient
$\mathscr F=\widetilde F$ un $\mathscr O_X$-Module quasi-cohérent,
$\mathscr G=\widetilde G$ un $\mathscr O_Y$-Module quasi-cohérent;
considérant $\mathscr F$ et $\mathscr G$ comme des complexes de Modules, on
notera $\mathscr{T}\!or_n^{\mathscr O_S}(\mathscr F,\mathscr G)$ ou
$\mathscr{T}\!or_n^S(\mathscr F,\mathscr G)$ leur hypertor d'indice $n$; il
résulte de (6.3.5 (ii)) que l'on a
\[
  \mathscr{T}\!or_n^{\mathscr O_S}(\mathscr F,\mathscr G)
  =\bigl(\operatorname{Tor}_n^A(F,G)\bigr)^{\sim}.
  \tag{6.4.3.1}\label{III.6.4.3.1.fr}
\]

Revenons alors au cas général de deux complexes de Modules quasi-cohérents
$\mathscr P_\bullet$, $\mathscr Q_\bullet$. Les formules (6.4.1.1) et
(6.4.3.1) montrent, compte tenu de la prop. (6.3.2), que
$\mathscr{T}\!or_\bullet^S(\mathscr P_\bullet,\mathscr Q_\bullet)$ est
l'aboutissement de deux suites spectrales
$'\mathscr E(\mathscr P_\bullet,\mathscr Q_\bullet)$,
$''\mathscr E(\mathscr P_\bullet,\mathscr Q_\bullet)$, dont les termes
$E_2$ sont donnés par
\[
  '\mathscr E_{pq}^2
  =\mathscr H_p\bigl(\mathscr{T}\!or_q^S
    (\mathscr P_\bullet,\mathscr Q_\bullet)\bigr)
  \tag{6.4.3.2}\label{III.6.4.3.2.fr}
\]
\[
  ''\mathscr E_{pq}^2
  =\bigoplus_{q'+q''=q}
  \mathscr{T}\!or_p^S
  \bigl(\mathscr H_{q'}(\mathscr P_\bullet),
        \mathscr H_{q''}(\mathscr Q_\bullet)\bigr)
  \tag{6.4.3.3}\label{III.6.4.3.3.fr}
\]
où $\mathscr{T}\!or_q^S(\mathscr P_\bullet,\mathscr Q_\bullet)$ est le
bicomplexe de $\mathscr O_{X\times_S Y}$-Modules quasi-cohérents
$\mathscr{T}\!or_q^S(\mathscr P_i,\mathscr Q_j)$.
\end{env}

\begin{env}[6.4.4]
Considérons maintenant deux autres schémas affines
$X^{(1)}=\operatorname{Spec}(B^{(1)})$,
$Y^{(1)}=\operatorname{Spec}(C^{(1)})$, où $B^{(1)}$ et $C^{(1)}$ sont des
$A$-algèbres, et supposons donnés deux

